%%%%%%%%%%%%%%%% Präambel-ANFANG
% Die folgenden Zeilen zwischen Präambel-ANFANG und Präambel-ENDE bitte nicht ändern.

% Das Latex-Dokument beginnt mit dem \documentclass-Befehl, welcher eine angepasste Version der ACM sig-alternate-Klasse (ewks-latex) referenziert. Die Datei ewks-latex.cls muss sich im gleichen Verzeichnis wie die Zusammenfassung befinden.
\documentclass[ngerman]{ewks-latex}

% Anpassungen an die deutsche Sprache.
\usepackage[ngerman]{babel}

% Angaben zu Zeichensätzen (unter anderem ist Nutzung deutscher Umlaute möglich).
\usepackage[T1]{fontenc}
\usepackage[utf8]{inputenc}

% Ohne dieses Package können keine eps Bilder angezeigt werden
\usepackage{epstopdf}

% Entfernung von für den Copyright-Hinweis der Konferenz reserviertem Platz in der ersten Spalte.
\usepackage{etoolbox}
\makeatletter
\patchcmd{\maketitle}{\@copyrightspace}{}{}{}
\makeatother

% Anzahl der Autoren der Zusammenfassung (in diesem Semester zwei Studenten).
\numberofauthors{2}

%%%%%%%%%%%%%%%% Präambel-ENDE

% Das lipsum-Paket wird für die Generierung von Blindtext in dieser Beispieldatei verwendet und kann aus der tatsächlichen Abgabe entfernt werden.
\usepackage{lipsum}

% Das listings-Paket kann zum Setzen von Quellcode verwendet werden. Es kann auch entfernt werden, falls kein Quellcode gesetzt werden muss.
\usepackage{listings}

% Start des Dokuments.
\begin{document}

% Angabe des *englischen* Titels des Papers.
% Bitte das Tag Approach to Define Highly Scalable Metamodels Based on JSON entsprechend ersetzen.
\title{Approach to Define Highly Scalable Metamodels Based on JSON}

% Als Untertitel sollte darauf hingewiesen werden, dass es sich um eine Zusammenfassung des Papers der jeweiligen Konferenz inkl. Jahrgang handelt.
% Bitte das Tag 2015 entsprechend ersetzen.
\subtitle{Zusammenfassung des auf der BigMDE 2015 veröffentlichten, gleichnamigen Papers}

% Angaben zu den Autoren der *Zusammenfassung*.
% Bitte die Tags <VORNAMEx>, <NAMEx>, <MATRIKELNRx>, <EMAILx> durch die eigenen Angaben ersetzen.
\author{
	% Erster Autor (i. d. R. der Corresponding Author, der auch die Mail zur Gruppenanmeldung und Papier-Auswahl verfasst hat).
	\alignauthor
	<Peter> <S>\\
	\affaddr{Fachhochschule Dortmund, Fachbereich Informatik}\\
    \affaddr{<MATRIKELNR1>}\\
    \email{ps@fh-dortmund.de}    
	% Zweiter Autor (diese Kommentarzeile nicht entfernen!)
	\alignauthor
	<Sascha> <D>\\
	\affaddr{Fachhochschule Dortmund, Fachbereich Informatik}\\
    \affaddr{<MATRIKELNR2>}\\
    \email{<EMAIL2>@stud.fh-dortmund.de}    
}

% Titel setzen. Dieser Befehl sollte an dieser Stelle belassen werden und vor dem eigentlichen Text auftauchen.
\maketitle

% Das Paper wie auch die Zusammenfassung gliedern sich in mehrere Abschnitte (\section-Befehl). Der erste Abschnitt sollte der Einleitung (Introduction) vorbehalten sein und eine Spalte nicht überschreiten.
\section{Einleitung}
Das Paper Approach to Define Highly Scalable Metamodels Based on JSON \cite{gerhart2015approach} wurde im Rahmen der BigMDE 2015 veröffentlicht und es beinhaltet \textit{einen Ansatz} zur Speicherung von \textbf{D}omain-\textbf{S}pecific \textbf{M}odeling \textbf{L}anguages (DSML). DSMLs werden in der Softwareentwicklung eingesetzt, um verschiedene Diagramme in textueller Form abzubilden. Dabei geht es nicht um eine gut lesbare und verständliche Übersicht, sondern um die Generierung der Diagramme. Hierbei gibt es verschiedene Ansätze zur Generierung. So können die Diagramme beispielsweise direkt durch interpretierten Java Quellcode, spezielle \textit{@Annotationen} oder eine zusätzliche Sprache generiert werden\cite{france2005domain}. In der Vergangenheit wurde dafür oft auf XML zurückgegriffen. Für Softwareprojekte, die mit wenigen Klassen und Referenzen auf der Programmierebene auskommen ist dies unproblematisch. In großen Projekten ist die Anzahl der beiden Elemente aber unter Umständen so hoch, dass die angelegten XML-Dateien mehrere Gigabyte Speicher beanspruchen. Hier kommt ein weiterer Nachteil des XML-Formats zum tragen. Modelle können nur als Ganzes ausgelesen werden. \\
Diese Limitierungen in der Skalierbarkeit können durch das JSON-Format verbessert und zum Teil sogar aufgehoben werden. Beispielsweise kann der Speicherbedarf durch einen geringeren Overhead im Quellcode reduziert werden und Modelle können über sogenannte \textit{Identifier} partiell geladen werden. \\
Um die zuvor getroffenen Aussagen zu untermauern und die Vorteile des JSON-Formats faktisch darzustellen wird im Folgenden eine JSON-DSML-Implementierung beschrieben und in einem anschließenden Benchmark mit Ecore verglichen. \\
\textit{Die Gliederung dieser Zusammenfassung orientiert sich am Originalpaper... \glqq So etwas ähnliches stand in beiden Positivbeispielen.\grqq} 
\section{Related Work}


\section{Hinweise zu den folgenden Abschnitten}
Nach der Einleitung sollte in den folgenden Abschnitten der Inhalt des Papers vorgestellt werden. Die Einteilung der Abschnitte kann sich dabei am Paper orientieren.

\section{Hinweise zur Formatierung}
Ein paar Beispiele zur Formatierung von Text:
\begin{enumerate}
\item \textit{kursiv}
\item \textbf{fett}
\item \texttt{Code}
\end{enumerate}

\section{Tabellen}
\begin{table} % Einleitung der Tabelle
\centering % Überschrift zentriert
\caption{Ein paar Java-Schlüsselwörter} % Überschrift
\label{tab:example} % Label für spätere Referenzierung per \ref
\begin{tabular}{|c|c|} \hline % Einleitung des Tabelleninhalts
Schlüsselwort & Bedeutung \\ \hline % Zeile 1; Spaltentrennung per &
\texttt{abstract} & Abstrakte Klasse/Methode \\ \hline % Zeile 2
\texttt{class} & Einleitung einer Klasse \\ \hline % Zeile 3
\texttt{throw} & Exception werfen \\ \hline % Zeile 4
\end{tabular}
\end{table}

Tabellen lassen sich mit Hilfe der \texttt{table}- und \texttt{tabular}-Umgebungen setzen. Als Beispiel für eine konkrete Anwendung siehe Tabelle \ref{tab:example}.

Warum ist meine Tabelle nicht dort, wo ich sie definiert habe? Weil die Vorlage den "`floating"'-Modus von \LaTeX verwendet und Tabellen, Abbildungen und andere Elemente so anordnet, dass wenig Leerzeilen entstehen.

Die Positionierung der Elemente sollte dabei nicht manuell geändert werden, außer es handelt sich um doppespaltige Tabellen, Abbildungen etc. Ein Beispiel für ein doppelspaltiges Element zeigt Abbildung \ref{fig:doppelspaltig}.

Doppelspaltige Elemente werden durch Varianten der einleitenden Befehle mit dem Suffix "`*"' gesetzt:
\begin{itemize}
\item Tabellen: \texttt{begin\{table*\}}
\item Abbildungen (siehe Abschnitt \ref{sec:figures}): \texttt{begin\{figure*\}}
\end{itemize}

\section{Abbildungen}
\label{sec:figures}
Abbildungen werden durch die \texttt{figure}-Umgebung gesetzt. Als Format empfiehlt sich PostScript (ps) oder Encapsulated PostScript (eps).

\begin{figure}
\centering
\epsfig{file=Quadrat.eps,height=1in,width=1in}
\caption{Ein einspaltiges graues Quadrat (eps-Format).}
\end{figure}

\begin{figure*}
\centering
\epsfig{file=Quadrat.eps,height=1in,width=1in}
\caption{Ein doppelspaltiges graues Quadrat (eps-Format).}
\label{fig:doppelspaltig}
\end{figure*}

\section{Verweise}
Mit Hilfe von \LaTeX \ lassen sich eine Reihe von Verweisen auf weiterführende Informationen nutzen.

Die wichtigste "`Verweistechnik"' ist der Literaturverweis. Der Befehl für einen Literaturverweis lautet \texttt{\textbackslash cite\{<Bib\TeX-MARKE>\}}. \texttt{<Bib\TeX-MARKE>} ist dabei durch die jeweilige Bib\TeX-Marke zu ersetzen. Bei Bib\TeX \ handelt es sich um ein Programm, mit dem sich Literaturverweise definieren lassen. Es bietet eine Anbindung an \LaTeX. Ein Beispiel für eine Bib\TeX-Datei ist "`literatur.bib"', die sich im selben Verzeichnis wie diese \texttt{.tex}-Datei befinden sollte. Der Befehl \texttt{\textbackslash cite\{latex\}} führt dann zu folgender Ausgabe: \cite{latex}.

Die Bib\TeX-Datei "`literatur.bib"' können Sie entweder mit einem Texteditor anpassen oder Sie nutzen ein Tool wie JabRef\footnote{\texttt{http://jabref.sourceforge.net}}. Fußnoten setzt man übrigens mit \texttt{\textbackslash footnote\{\}}.

\section{Formeln}
Eine große Stärke von \LaTeX \ ist das Setzen von Formeln. Formeln können \textit{inline} gesetzt werden, d. h. sie erscheinen im Fließtext: $\lim_{n\rightarrow \infty}x=0$.

Formeln werden entweder in zwei \$...\$ eingeschlossen (Kurzform) oder sind von \texttt{\textbackslash begin\{$x$\}}...\texttt{\textbackslash end\{$x$\}} umgeben mit $x\in\left\{ \mathtt{math\textrm{,}\:\mathtt{equation\textrm{,}\: displaymath}}\right\}$ (Langform).

Formeln können auch in abgesetzter Form erscheinen, undzwar...
\begin{itemize}
\item ... nummeriert (\texttt{equation}-Umgebung):
\begin{equation}
{n+1\choose k} = {n\choose k} + {n \choose k-1}
\end{equation}

\item ... nicht-nummeriert (\texttt{displaymath}-Umgebung):
\begin{displaymath}
|x| = \left\{ \begin{array}{rl}
 -x &\mbox{ falls $x<0$} \\
  x &\mbox{ sonst}
       \end{array} \right.
\end{displaymath}
\end{itemize} 



\section{Quellcode}
Ab und an muss auch Quelltext gesetzt werden. Dies geschieht mit Hilfe des \texttt{listings}-Pakets\footnote{\texttt{http://texdoc.net/texmf-dist/doc/latex/listings/\\listings.pdf}}. Listing \ref{lst:example} zeigt ein Beispiel für ein "`floating"'-Listing.

\begin{lstlisting}[float, caption=Beispiel eines Java-Listings, label=lst:example, language=java, basicstyle=\small, tabsize=4, showstringspaces=false, numbers=left, frame=single, keywordstyle=\bfseries]
public class Main {
	public static void main(String[] args) {
		System.out.println("Hallo Welt");
	}
}
\end{lstlisting}

\section{Abschnitt 8 (Blindtext)}
\lipsum[5-10]

\section{Abschnitt 9 (Blindtext)}
\lipsum[11-20]

\section{Abschnitt 10 (Blindtext)}
\lipsum[21-30]

\section{Abschnitt 11 (Blindtext)}
\lipsum[31-32]

\section{Related Work (Blindtext)}
\lipsum[33-37]

\section{Zusammenfassung (Blindtext)}
\lipsum[42-44]

%%%%%%%%%%%%%%%% Epilog-ANFANG
% Die folgenden Zeilen zwischen Epilog-ANFANG und Epilog-ENDE bitte nicht ändern.

% Neue Seite für das Literaturverzeichnis
\newpage

% Der voreingestellte BibTeX-Stil ist abbrv und sollte nicht geändert werden. Er führt dazu, dass die Quellenverweise durchnummeriert werden, was die Regel für wissenschaftliche Veröffentlichungen in den Naturwissenschaften ist.
\bibliographystyle{abbrv}
%%%%%%%%%%%%%%%% Epilog-ENDE

% Verweis auf die BibTeX-Datei. Diese sollte im gleichen Ordner wie die tex-Datei liegen. ACHTUNG: Der Dateiname ist ohne die Endung .bib anzugeben.
\bibliography{literatur}
\end{document}